% !TeX root = ../main.tex
% !TeX encoding = UTF-8

\chapter{The Compliance Clock, What the Deferral Did Not Defer}

The audit committee quarterly AI governance review. The General Counsel has just received a formal information request from the FCA. A specific AI-assisted lending decision made fourteen months ago is under examination. The regulator wants the Decision Identity, the Information Picture, and the Accountability Chain for that decision, not the current model documentation, not the governance framework, but the contemporaneous record of what existed at the moment that specific decision was made. Around the table, the question is the same one that has been asked in six industries across seven decades of this book's case studies: can we answer this?

The EU AI Act set 2 August 2026 as the date from which operators of high-risk AI systems would be bound by its core obligations. That date did not hold. The Digital Omnibus on AI (Regulation (EU) 2026/1744, in force since 27 July 2026) deferred the obligations for Annex III systems, which include credit scoring and recruitment, to 2 December 2027, and for systems embedded in regulated products under Annex I to 2 August 2028. The deferral was granted because the harmonised standards on which compliance depends were not ready, not because the obligations themselves changed. The Article 50 transparency duties were not deferred: they applied from 2 August 2026 as scheduled and are now enforceable. The penalty structure was untouched by the Omnibus and remains concrete: up to €15 million or 3\% of global annual turnover for operators of non-compliant high-risk systems. For a FTSE 350 company with annual turnover of £2 billion, 3\% of global annual turnover is £60 million. But the penalty is the consequence of failing to close the gap, and the gap is closed by the board, not by the calendar. This chapter is about using the time the deferral has bought to close it before the question arrives.

\begin{quote}
\emph{The deferral to 2 December 2027 is not a reprieve. It is a test. The obligations did not change; only the date on which their absence becomes penalisable did. The UK duties this book has examined, the SM\&CR evidential burden and the Companies Act duty of care, carried no deferral at all. A board that moves this file down the agenda because a date moved in Brussels is exhibiting precisely the deferral pattern this book was written to end.}
\end{quote}

\scenebreak

\section{The Concurrent Risk Problem, One Gap, Many Obligations}

The Digital Alibi gap is not a single regulatory risk. It is a concurrent risk across multiple frameworks simultaneously. The EU AI Act is the most visible obligation, but it is not the only one that the gap exposes. Each framework in the regulatory stack that applies to a regulated organisation makes its own demands on specific facets of the governance address. Understanding which framework is the constraining obligation on which facet is the first act of an effective governance investment strategy.

The table below maps the major AI governance standards and regulations to their minimum governance address requirements. Two observations are essential before reading it. First, the standards do not require identical addresses, they impose different minimum levels on different facets, and the most demanding requirement on each facet across the applicable standards is the binding floor for an organisation subject to all of them. Second, the final row of the table is the most important. An organisation whose AI governance programme is built around voluntary principles and industry commitments has not met any formal standard. It has documented its aspirations.

\begin{longtable}[]{@{}
  >{\raggedright\arraybackslash}p{(\columnwidth - 4\tabcolsep) * \real{0.4296}}
  >{\raggedright\arraybackslash}p{(\columnwidth - 4\tabcolsep) * \real{0.1408}}
  >{\raggedright\arraybackslash}p{(\columnwidth - 4\tabcolsep) * \real{0.4296}}@{}}
\toprule\noalign{}
\begin{minipage}[b]{\linewidth}\raggedright
\textbf{Standard}
\end{minipage} & \begin{minipage}[b]{\linewidth}\raggedright
\textbf{Minimum address}
\end{minipage} & \begin{minipage}[b]{\linewidth}\raggedright
\textbf{Pressure facet}
\end{minipage} \\
\midrule\noalign{}
\endhead
\bottomrule\noalign{}
\endlastfoot
EU AI Act, High Risk AI Systems (Arts. 9, 12, 13, 14, 17, 19) & \textbf{A4·E4·C4·R4·M4} & All five, the only standard requiring level 4 simultaneously \\
EU AI Act, General Purpose AI (Arts. 53-55) & \textbf{A3·E3·C3·R3·M3} & Convergence point, baseline obligations \\
EU AI Act, Prohibited AI Practices (Art. 5) & \textbf{A4·E4·C3·R4·M3} & A and R, identification and accountability \\
DORA 2022, ICT risk, resilience, third-party AI & \textbf{A3·E4·C4·R4·M3} & E and C, inventory completeness and verified controls \\
FCA SM\&CR, Named accountability, conduct rules & \textbf{A4·E3·C3·R3·M3} & A, named individual accountability is the sole pressure facet \\
GDPR, Article 22, automated decision-making & \textbf{A3·E3·C3·R3·M3} & Convergence point, minimum automated decision obligations \\
NIS2 2024, Critical infrastructure, supply chain & \textbf{A3·E4·C4·R4·M3} & E and C, mirrors DORA for critical infrastructure operators \\
ISO 42001:2023, AI management system & \textbf{A3·E4·C4·R3·M4} & E and C, inventory and independently verified controls \\
ISO 27001:2022, ISMS, certification & \textbf{A3·E3·C4·R4·M4} & C and M, control maturity and continuous improvement \\
NIST AI RMF, GOVERN, MAP, MEASURE, MANAGE & \textbf{A3·E3·C4·R3·M4} & C and M, operational governance maturity \\
UK AI Safety Framework, Safety by design & \textbf{A4·E3·C3·R3·M3} & A, contestability requirement drives A4 \\
OECD AI Principles / UNESCO / EU AI Pact & \textbf{A2·E2·C1-2·R2·M2} & Voluntary, no enforcement weight, no legal protection \\
\end{longtable}

The most important observation from this table is in the final row. An organisation that has signed the EU AI Pact, adopted the OECD AI Principles, and referenced the UNESCO Recommendation in its AI strategy documents is at approximately A2·E2·C1-2·R2·M2. Chapter 14 established that Level 2 is an exposure position with a principles overlay, it provides no legal protection and no defensible position under regulatory scrutiny. When an incident occurs, those voluntary commitments will be adduced as evidence that the organisation knew the standard it had not met.

For a regulated organisation subject to the EU AI Act High Risk classification, FCA SM\&CR, and DORA simultaneously, the binding floor on each facet is: A4 (driven by SM\&CR and EU AI Act High Risk equally), E4 (driven by DORA and EU AI Act High Risk), C4 (driven by DORA, EU AI Act High Risk, and ISO 42001), R4 (driven by EU AI Act High Risk and DORA), M4 (driven by EU AI Act High Risk). The combined requirement is A4·E4·C4·R4·M4, the full sector ceiling for financial services and insurance identified in Chapter 14. Every facet at level 4. The Digital Alibi Assessment that Chapter 16 described produces the confirmed address from which the distance to that ceiling can be measured.

\scenebreak

\section{The A4 Problem, Why Accountability Is the Hardest Facet}

The 2 December 2027 obligation is not evenly distributed across the five facets. The technical implementation work, model versioning, data lineage capture, parameter logging, oversight recording, is demanding but has established tooling and defined timelines. The regulatory mapping, identifying which systems fall within the EU AI Act's high-risk classification and documenting the obligations that apply, is time-consuming but procedurally clear. The hardest facet to advance is A. Not because the requirement is unclear. Because reaching A4 requires a change to how a board understands its own accountability that most boards have not yet made.

A4 requires formal accountability chains documented at board level. Not a named CISO. Not a Chief AI Officer in the executive team. A board-documented ownership structure in which a named individual at board level has accepted the specific risk that each material AI system represents, documented contemporaneously, updated at each material change to the system or the accountability structure, and retrievable with forensic precision for any specific decision moment. The Accountability Chain field in the Decision Map from Chapter 17 is the operational implementation of A4. Before that field can be populated for any decision, the named individual must exist and their acceptance must be documented.

The most common failure mode on the path to A4 is the interim appointment gap. An organisation that has named its Chief Risk Officer as the A4 accountability holder for AI governance finds, when it examines the Decision Map for a specific decision made during a period when the CRO role was vacant between appointments, that the chain has no anchor at the board level for that decision period. The gap is not a failure of governance intent. It is a failure of governance continuity, the same failure the Lloyd's Names catastrophe produced at a different scale. The A4 requirement is met continuously, not periodically.

\begin{quote}
\emph{Reaching A4 is not a governance programme. It is a governance commitment, to named individual accountability, documented contemporaneously, continuous across every material AI decision moment including the ones made between appointments, during restructures, and across reporting lines that change. Most organisations underestimate how far they are from that commitment.}
\end{quote}

\scenebreak

\section{The 90-Day Protocol, A Prioritised Remediation Sequence}

For an organisation whose board has concluded, after reading this book, that it needs to close its Digital Alibi gap without waiting for the deferred date, the following 90-day sequence is the prioritised governance response. It is not a comprehensive AI governance programme. It is the minimum sequence of actions required to move from an unconfirmed governance address to a confirmed one, and from a confirmed address to the governance architecture that closes the gap on the facets most immediately exposed, under the UK duties that already apply and the EU obligations that apply from 2 December 2027.

\section{Days 1-15 AI Systems Inventory}

Identify every AI system contributing to high-risk decisions across the five material decision types established in Chapter 16. The inventory must include production systems, vendor-embedded AI in licensed platforms, and shadow AI deployed without formal governance approval. Map each identified system to the EU AI Act's Article 9 risk categories to determine which fall within the high-risk classification and are therefore subject to the high-risk obligations that apply from 2 December 2027.

The E facet is consistently the most underestimated at self-assessment. Most organisations that believe they are at E3 discover, when vendor-embedded AI is properly surfaced, that they are at E2. The inventory exercise corrects that underestimate. It should be conducted by the function responsible for the formal AI inventory working with the business units that operate the material decision processes, not by asking the technology team to confirm the accuracy of the existing inventory, which will produce a self-attestation rather than a confirmation.

The inventory must also address the Vanishing Vendor risk within the compliance window. If a vendor providing a material AI system enters administration, is acquired, or discontinues the product before the organisation has built its Decision Map, the historical decisions that system contributed to become permanently irrecoverable on the Decision Identity and Accountability Chain fields. Chapter 11 established the Vanishing Vendor as the accountability failure mode for agentic AI deployments. In the context of the compliance obligations, it applies to every vendor-provided AI system in the material decision inventory: if the organisation does not hold the model documentation, version history, and parameter configuration independently, separate from the vendor's infrastructure, then vendor insolvency before Day 1 of the 90-day protocol produces exactly the governance gap that the Protocol exists to close. The inventory exercise must confirm, for each vendor-provided system, that the organisation holds the documentation it would need to populate the Decision Map if the vendor ceased operating tomorrow.

\section{Days 16-30 Digital Alibi Gap Assessment}

Conduct the Digital Alibi Assessment described in Chapter 16 across the five most material AI-assisted decisions in the inventory. Produce the Decision Reconstruction Report. Confirm the governance address across A·E·C·R·M with forensic precision. The confirmed address is the starting point for every subsequent investment decision.

The gap between the self-assessed address and the confirmed address is itself a governance finding that must be presented to the audit committee or board risk committee. If the gap is significant, if the confirmed address is two or more levels below the self-assessed address on any facet, it must be presented to the full board. The board cannot govern the gap it does not know exists.

\section{Days 31-50 Governance Architecture Design}

Design the Decision Map architecture for each material AI system in the confirmed inventory. Name individual accountabilities to board level, meeting the A4 requirement. The accountability nomination must cover every material AI system, must name a board-level risk acceptance authority for each, and must address the continuity requirement, the governance architecture that ensures the chain does not break during interim appointments or role transitions.

Address the negative space for each decision type, the data absences and design exclusions that the Negative Space Record field will need to capture. The negative space documentation is the field most commonly approached last and most consistently the one that produces the largest evidential gap when challenged. It should be addressed during architecture design, not during technical implementation.

\section{Days 51-70 Technical Implementation}

Implement model versioning, parameter logging, data lineage capture, and human oversight recording at output-moment resolution for each material AI system. The implementation must produce a complete Decision Map entry, all six fields, for each governed decision, automatically, without manual assembly. An implementation that captures five of the six fields and requires manual assembly for the sixth is not complete.

Design and validate the Alibi Record architecture, the tamper-evident, contemporaneous, actively retrievable governance record that Chapter 20 describes as the destination of the Protocol. The Alibi Record cannot be built on a governance address below A3·E3·C3·R3·M3. The technical implementation phase confirms that the prerequisite address has been reached before attempting to implement the Alibi Record.

\section{Days 71-85 Board Briefing and Sign-Off}

Present the confirmed governance address, the Decision Reconstruction Report, and the governance architecture design to the audit committee or board risk committee. Obtain formal board sign-off on the accountability chains, the named individuals at board level who are accepting accountability for each material AI system. This is the governance act that advances the address from A3 to A4. Without named sign-off documented at board level, the accountability chain does not exist at the level the EU AI Act and SM\&CR require.

The board sign-off is not a resolution noting the governance programme. It is a record of named individuals accepting specific accountabilities for specific AI systems at a specific date. It must be minuted with that specificity. A resolution that notes the governance programme without naming individuals and systems has not produced A4.

\section{Days 86-90 Documentation Retention Architecture}

Confirm that all Decision Map entries and Alibi Records are stored in a form that satisfies the EU AI Act's Article 12 record-keeping obligations, tamper-evident, retrievable on demand, retained for the period the Act requires. The storage architecture must survive the tests established in the Retrievability Certification from Chapter 17: complete retrieval, by an independent examiner, within the required timeframe.

Review the escalation protocol for capture failures. Confirm that the named Senior Manager on the Accountability Chain for each material AI system will receive direct notification of any Decision Map capture failure within the session in which it occurs. Test the escalation mechanism before it is needed, not after.

\begin{longtable}[]{@{}
  >{\raggedright\arraybackslash}p{(\columnwidth - 0\tabcolsep) * \real{1.0000}}@{}}
\toprule\noalign{}
\begin{minipage}[b]{\linewidth}\raggedright
\textbf{WHERE MOST ORGANISATIONS STAND TODAY, AND WHAT THAT MEANS}
\end{minipage} \\
\midrule\noalign{}
\endhead
\bottomrule\noalign{}
\endlastfoot
At A2·E2·C2·R2·M2 or below: no formal compliance standard met. The high-risk obligations that apply from 2 December 2027 cannot be met from this address without immediate investment, and the deferral suspended none of the UK exposure. SM\&CR personal liability unaddressed. The Digital Alibi gap is total, no decision in any material AI system is defensible from the contemporaneous record. \\
At A3·E3·C3·R3·M3 (convergence point): baseline formal compliance obligations engaged. EU AI Act High Risk requirements not fully met, A4 and E4 are still outstanding. DORA C4 and E4 requirements not met. The Digital Alibi gap is partial, some decisions have partial contemporaneous records but none have a complete Decision Map entry. \\
At A4·E4·C4·R4·M4: EU AI Act High Risk compliance defensible. DORA and SM\&CR obligations met. Decision Map architecture operational. The Digital Alibi gap is closed for governed decisions, the board can answer the output-moment question for every decision in the governed inventory from the contemporaneous record. \\
The distance between where most regulated organisations currently sit (A2·E2·C2·R2·M2 to A3·E3·C3·R3·M3, based on confirmed address distributions from advisory practice) and where the obligation will require them to be (A4·E4·C4·R4·M4) is the governance investment the 90-day Protocol is designed to close. It is not achievable in 90 days for every organisation. But it is achievable for every organisation that starts on Day 1. \\
\end{longtable}

\scenebreak

\section{After the Deadline, What the Governance Address Must Sustain}

Reaching A4·E4·C4·R4·M4 is the objective the 2 December 2027 date sets; for an organisation whose confirmed address sits below it, the distance is measured in the time the deferral has left and in the UK exposure that already applies. Sustaining it is the ongoing objective. Chapter 14 established the concept of governance address decay, the movement of an address down over time as systems change, personnel move, regulations evolve, and the governance architecture fails to keep pace. A4 sustained requires that the named individual accountabilities are updated within thirty days of any change to the accountability structure. E4 sustained requires that the AI systems inventory is reviewed quarterly and updated within thirty days of any new system deployment or vendor change. C4 sustained requires that the independent verification cycle is maintained and that any control failure identified in that cycle is escalated and remediated within a defined timeframe.

The governance address is not a certification. It is a position. Positions move. An organisation that reached A4·E4·C4·R4·M4 in December 2027 and has not maintained the governance architecture that produced that address will find, twelve months later, that its confirmed address has decayed, and that the decisions made during the decay period are not defensible from the contemporaneous record. The Governance Retainer that Chapter 21 describes is the mechanism for preventing that decay.

Chapter 19 addresses the board conversation that must happen to sustain the governance address, the communication architecture for a board that has received its first Decision Reconstruction Report, understands its confirmed address, and must now make the governance investments the Protocol requires without losing momentum or producing organisational paralysis. The 90-day Protocol closes the gap. The board conversation sustains the closure.
